% ============================================================
% TUTORIAL: Fase 5 — RAG con OCI (Oracle Cloud Infrastructure)
% Proyecto: Sentinel AcademIA
% Compilar: tectonic tutorial-fase5-rag.tex
% ============================================================

\documentclass[11pt,a4paper,oneside]{article}

\usepackage[utf8]{inputenc}
\usepackage[T1]{fontenc}
\usepackage[spanish,es-nodecimaldot]{babel}
\usepackage{lmodern}
\usepackage[margin=2.5cm]{geometry}
\usepackage{parskip}
\usepackage{microtype}

\usepackage{xcolor}
\usepackage{colortbl}
\usepackage{array}
\usepackage{booktabs}
\usepackage{amssymb}
\usepackage{pifont}
\usepackage{tcolorbox}
\tcbuselibrary{breakable,skins}

\usepackage{listings}

\lstdefinelanguage{yaml}{
  basicstyle=\ttfamily\small,
  keywordstyle=\color{blue!70!black}\bfseries,
  stringstyle=\color{red!70!black},
  commentstyle=\color{green!50!black}\itshape,
  morekeywords={Type,Properties,Version,Statement,Effect,Principal,Action,Resource},
}

\lstdefinelanguage{bash}{
  basicstyle=\ttfamily\small,
  keywordstyle=\color{blue!70!black}\bfseries,
  stringstyle=\color{red!70!black},
  commentstyle=\color{green!50!black}\itshape,
  morekeywords={oci,aws,npm,cd,ls,mkdir,rm,cp,mv,export,unset,echo,cat,curl,sleep,npx,git,python3},
}

\lstdefinelanguage{json}{
  basicstyle=\ttfamily\small,
  showstringspaces=false,
  breaklines=true,
  literate=
    *{\{}{{{\{}}}{1}
     {\}}{{{\}}}}{1}
     {[}{{{[}}}{1}
     {]}{{{]}}}{1}
     {:}{{{:}}}{1}
     {,}{{{,}}}{1},
  string=[s]{"}{"},
  stringstyle=\color{red!70!black},
  morestring=[b]",
}

\lstset{
  basicstyle=\ttfamily\small,
  numberstyle=\tiny\color{gray},
  numbers=left,
  numbersep=8pt,
  frame=single,
  framerule=0.4pt,
  rulecolor=\color{gray!50},
  backgroundcolor=\color{gray!5},
  breaklines=true,
  breakatwhitespace=true,
  showstringspaces=false,
  captionpos=b,
  tabsize=2,
  literate={á}{{\'a}}1 {é}{{\'e}}1 {í}{{\'i}}1 {ó}{{\'o}}1 {ú}{{\'u}}1
           {ñ}{{\~n}}1 {Á}{{\'A}}1 {É}{{\'E}}1 {Í}{{\'I}}1 {Ó}{{\'O}}1 {Ú}{{\'U}}1
           {¿}{{?`}}1 {¡}{{!`}}1
}

\usepackage[colorlinks=true,linkcolor=blue!60!black,urlcolor=blue!60!black,citecolor=blue!60!black]{hyperref}

\usepackage{fancyhdr}
\pagestyle{fancy}
\fancyhf{}
\fancyhead[L]{\small\textsc{Sentinel AcademIA}}
\fancyhead[R]{\small Tutorial Fase 5}
\fancyfoot[C]{\thepage}
\renewcommand{\headrulewidth}{0.4pt}

\definecolor{okgreen}{RGB}{40,140,60}
\definecolor{errred}{RGB}{200,50,50}
\definecolor{warnorange}{RGB}{220,130,30}
\definecolor{infoblue}{RGB}{30,90,180}

\newcommand{\ok}{\textcolor{okgreen}{\textbf{\checkmark}}}
\newcommand{\err}{\textcolor{errred}{\textbf{\ding{55}}}}
\newcommand{\warn}{\textcolor{warnorange}{\textbf{\textasciitilde}}}

\title{\Huge\textbf{Fase 5: RAG con OCI} \\[0.3em]
       \Large Retrieval Augmented Generation con Oracle Cloud \\[0.5em]
       \normalsize Tutorial paso a paso}
\author{Proyecto Sentinel AcademIA \\ \small lFunknown}
\date{\today}

\begin{document}

\maketitle
\thispagestyle{empty}

\begin{tcolorbox}[colback=blue!5!white,colframe=infoblue,title={\textbf{¿Qué construimos en este tutorial?}}]
En la \textbf{Fase 5} agregamos un \textbf{chatbot con RAG} (Retrieval Augmented Generation) que responde preguntas sobre el reglamento universitario bas\'andose en el documento oficial.

\begin{itemize}
  \item \textbf{Knowledge base}: \texttt{reglamento.txt} subido a OCI Object Storage
  \item \textbf{Retrieval}: keyword search (RAG h\'ibrido) o OCI Generative AI Agent (producci\'on)
  \item \textbf{LLM}: Cohere Command (real) o MockLLMClient (Lambda)
  \item \textbf{Endpoint}: \texttt{POST /api/chat} con answer + sources (citations)
\end{itemize}

\textbf{Decisiones clave}:
\begin{itemize}
  \item \textbf{Por qu\'e OCI}: Oracle tiene Generative AI Agents managed, ideal para RAG
  \item \textbf{Por qu\'e keyword search en Lambda}: el package \texttt{oci} pesa 229MB (excede 250MB de Lambda)
  \item \textbf{Demo vs Producci\'on}: en prod se usa OCI Agent; en demo se hace keyword search en Python
\end{itemize}
\end{tcolorbox}

\begin{tcolorbox}[colback=warnorange!5!white,colframe=warnorange,title={\textbf{Herramientas de Oracle usadas}}]
\begin{itemize}
  \item \textbf{OCI Object Storage}: bucket para documents de la knowledge base
  \item \textbf{OCI Python SDK (v2.179.0)}: para crear bucket y subir archivos
  \item \textbf{OCI Generative AI Agents} (documentado, no desplegado): para RAG production-ready
  \item \textbf{Oracle CLI 3.87.0}: para inspecci\'on de recursos
  \item \textbf{Python 3.14.3 del OCI installer}: \texttt{\textasciitilde/lib/oracle-cli/bin/python3}
\end{itemize}
\end{tcolorbox}

\tableofcontents
\newpage

% ============================================================
\section{¿Qu\'e es RAG?}
% ============================================================

\subsection{Definici\'on}

\textbf{RAG (Retrieval Augmented Generation)} = patr\'on donde:
\begin{enumerate}
  \item \textbf{Usuario} hace una pregunta
  \item \textbf{Retriever} busca documentos relevantes en una knowledge base
  \item \textbf{Contexto relevante} se inyecta en el prompt del LLM
  \item \textbf{LLM} genera una respuesta basada en el contexto (no en su conocimiento previo)
\end{enumerate}

\subsection{¿Por qu\'e RAG?}

\begin{itemize}
  \item \ok \textbf{Conocimiento actualizado}: el LLM no sabe del reglamento, pero RAG s\'i
  \item \ok \textbf{Citations}: podemos mostrar "esto viene del art\'iculo X"
  \item \ok \textbf{Menos alucinaciones}: el LLM responde bas\'andose en docs reales
  \item \ok \textbf{Dominio espec\'ifico}: el LLM general + el conocimiento espec\'ifico de tu org
\end{itemize}

\subsection{Diferencia con fine-tuning}

\begin{center}
\begin{tabular}{lll}
\toprule
& \textbf{Fine-tuning} & \textbf{RAG} \\
\midrule
\textbf{Costo} & Alto (entrenar el modelo) & Bajo (solo embeddings) \\
\textbf{Actualizaci\'on} & Re-entrenar & Solo actualizar docs \\
\textbf{Citations} & No & S\'i \\
\textbf{Hallucinations} & Posibles & Menos probables \\
\textbf{Setup} & Semanas & Horas \\
\bottomrule
\end{tabular}
\end{center}

% ============================================================
\section{Arquitectura del RAG}
% ============================================================

\subsection{Diagrama del flujo}

\begin{center}
\begin{tabular}{ccccccc}
\textbf{Cliente} & $\rightarrow$ & \textbf{API} & $\rightarrow$ & \textbf{knowledge\_service} & $\rightarrow$ & \textbf{Contexto} \\
"Que hago si & & POST /api/chat & & keyword search & & articulos \\
sufro acoso?" & & & & & & relevantes \\
\end{tabular}
\end{center}

\begin{center}
\begin{tabular}{ccccc}
\textbf{Contexto} & $\rightarrow$ & \textbf{LLM} & $\rightarrow$ & \textbf{Respuesta} \\
+ pregunta & & Cohere / Mock & & con citations \\
\end{tabular}
\end{center}

\subsection{Decisiones de dise\~no}

\begin{itemize}
  \item \ok \textbf{Knowledge base}: 8 documentos curados del reglamento
  \item \ok \textbf{Retrieval}: tokenizaci\'on + keyword match (no embeddings)
  \item \ok \textbf{Top-K}: 3 documentos m\'as relevantes
  \item \ok \textbf{LLM}: Cohere con system prompt estricto
  \item \ok \textbf{Citations}: cada source incluye id, title, content, score
\end{itemize}

% ============================================================
\section{Herramientas de Oracle usadas}
% ============================================================

\subsection{OCI Object Storage}

\textbf{Qu\'e es}: servicio de almacenamiento de objetos (como S3 de AWS) en Oracle Cloud.

\textbf{Caracter\'isticas}:
\begin{itemize}
  \item \ok Hasta \textbf{10 TB por objeto} y \textbf{ilimitado} en total
  \item \ok \textbf{Durabilidad 99.999999999\%} (11 nueves)
  \item \ok Compatible con \textbf{S3 API} (puedes usar boto3)
  \item \ok \textbf{Tier Standard} (frecuente) o \textbf{Archive} (raro)
  \item \ok Encryption at rest por default
\end{itemize}

\subsection{OCI Python SDK}

\textbf{Qu\'e es}: biblioteca oficial de Oracle para Python.

\begin{lstlisting}[language=bash,caption={Instalaci\'on}]
pip install oci
# Versi\'on actual: 2.179.0
\end{lstlisting}

\begin{lstlisting}[language=Python,caption={Uso b\'asico}]
import oci
from oci.config import from_file

# Cargar config desde ~/.oci/config
config = from_file("~/.oci/config", "DEFAULT")

# Crear cliente de Object Storage
object_storage = oci.object_storage.ObjectStorageClient(config)

# Listar buckets
buckets = object_storage.list_buckets(
    namespace_name="ax59x2so1pxw",
    compartment_id="ocid1.compartment.oc1..aaa..."
)
for bucket in buckets.data:
    print(f"  - {bucket.name}")
\end{lstlisting}

\textbf{Conceptos clave}:
\begin{itemize}
  \item \ok \textbf{Namespace}: string \'unico del tenancy (no OCID). Ej: \texttt{ax59x2so1pxw}
  \item \ok \textbf{Compartment ID}: OCID donde se crea el bucket
  \item \ok \textbf{Region}: donde est\'a el bucket. Ej: \texttt{us-chicago-1}
\end{itemize}

\begin{tcolorbox}[colback=errred!5!white,colframe=errred,title={Trampa \#1: namespace vs OCID}]
\textbf{Cu\'ando}: intentas crear bucket. \\
\textbf{Causa}: pasamos el OCID del tenancy como namespace. \\
\textbf{Soluci\'on}: el namespace es un string corto (\texttt{ax59x2so1pxw}), no el OCID.

\begin{lstlisting}
# MAL
namespace = "ocid1.tenancy.oc1..aaa..."  # 90+ caracteres

# BIEN
namespace = "ax59x2so1pxw"  # nombre corto del namespace
\end{lstlisting}
\end{tcolorbox}

\subsection{OCI Generative AI Agents (producci\'on)}

\textbf{Qu\'e es}: servicio managed de RAG de Oracle.

\textbf{Componentes}:
\begin{itemize}
  \item \ok \textbf{Data Source}: apunta a Object Storage bucket o folder
  \item \textbf{Knowledge Base}: agrupa data sources, configura chunking
  \item \textbf{Index}: vector index para b\'usqueda sem\'antica
  \item \textbf{Agent}: usa la knowledge base + LLM
  \item \textbf{Endpoint}: HTTP endpoint para invocar el agent
\end{itemize}

\textbf{Ventaja}: Oracle maneja TODO el pipeline (chunking, embeddings, vector store, retrieval, LLM call, citations).

\textbf{Limitaci\'on para nosotros}: el SDK pesa 229MB (excede 250MB de Lambda).

% ============================================================
\section{Configuraci\'on de OCI}
% ============================================================

\subsection{Verificar la configuraci\'on existente}

Ya ten\'iamos (de Fase 0) el archivo \texttt{\textasciitilde/.oci/config}:

\begin{lstlisting}[language=bash,caption={\textasciitilde/.oci/config}]
[DEFAULT]
user=ocid1.user.oc1..aaaaaaaa35wm3olju2qhtck4c75fq6xlitzugos7bygbnkebr4nucrwymdfq
fingerprint=33:c3:5e:8d:b7:d6:73:83:5e:78:7a:8d:6a:41:10:a6
tenancy=ocid1.tenancy.oc1..aaaaaaaaltcvasuc6yawzyqprgk6uvbac3kghmwg3ayqs5amh4ircglalzwa
compartment-id=ocid1.compartment.oc1..aaaaaaaa4vvf5a7lznbrfxporgxgmkr3mdanyktzsaosbwqc2urrq4pwwmzq
region=us-chicago-1
key_file=/home/lFunknown/.oci/private_key.pem
\end{lstlisting}

\subsection{Verificar la conexi\'on}

\begin{lstlisting}[language=bash,caption={Test de conexi\'on OCI}]
~/lib/oracle-cli/bin/python3 << 'PYEOF'
import oci
from oci.config import from_file

config = from_file("~/.oci/config", "DEFAULT")
identity = oci.identity.IdentityClient(config)

tenancy = identity.get_tenancy(config["tenancy"]).data
print(f"User: {config['user']}")
print(f"Tenancy: {tenancy.name}")
print(f"Region: {config['region']}")
PYEOF
\end{lstlisting}

\textbf{Output esperado}:
\begin{lstlisting}
User: ocid1.user.oc1..aaaaaaaa35wm3olju2qhtck4c75fq6xlitzugos7bygbnkebr4nucrwymdfq
Tenancy: fabricioladera
Region: us-chicago-1
\end{lstlisting}

% ============================================================
\section{Paso 1: Crear el OCI Object Storage bucket}
% =========================================================%

\subsection{Encontrar el namespace correcto}

El \textbf{namespace} de Object Storage NO es el OCID del tenancy, es un string corto. Para encontrarlo:

\begin{lstlisting}[language=bash,caption={Obtener namespace}]
~/lib/oracle-cli/bin/python3 << 'PYEOF'
import oci
from oci.config import from_file

config = from_file("~/.oci/config", "DEFAULT")
identity = oci.identity.IdentityClient(config)
namespace = identity.get_tenancy(config["tenancy"]).data
print(f"Namespace: {namespace.id}")
PYEOF
# Output: Namespace: ocid1.tenancy.oc1..aaaaaaaaltcvasuc6yawz...
\end{lstlisting}

Wait, eso es el OCID. Necesitamos el \textbf{Object Storage namespace}, que es diferente:

\begin{lstlisting}[language=bash,caption={Obtener Object Storage namespace}]
oci os ns get
# Output: {
#   "data": "ax59x2so1pxw"
# }
\end{lstlisting}

\subsection{Crear el bucket}

\begin{lstlisting}[language=Python,caption={Script para crear bucket}]
import oci
from oci.config import from_file

config = from_file("~/.oci/config", "DEFAULT")
object_storage = oci.object_storage.ObjectStorageClient(config)
compartment_id = "ocid1.compartment.oc1..aaaaaaaa4vvf5a7lznbrfxporgxgmkr3mdanyktzsaosbwqc2urrq4pwwmzq"

bucket_name = "sentinel-knowledge"
namespace = "ax59x2so1pxw"  # <-- el namespace real, NO el OCID

details = oci.object_storage.models.CreateBucketDetails(
    name=bucket_name,
    compartment_id=compartment_id,
    public_access_type="NoPublicAccess",
    storage_tier="Standard",
    versioning="Disabled",
)

try:
    result = object_storage.create_bucket(namespace, details)
    print(f"Bucket {bucket_name} creado!")
except oci.exceptions.ServiceError as e:
    if e.status == 400:
        print(f"Namespace mismatch: {e.message}")
    elif e.status == 409:
        print(f"Bucket ya existe")
    else:
        raise
\end{lstlisting}

\textbf{Output}:
\begin{lstlisting}
Bucket sentinel-knowledge creado!
  Created: 2026-06-20 15:07:11.973000+00:00
\end{lstlisting}

% ============================================================
\section{Paso 2: Crear el documento del reglamento}
% =========================================================%

\subsection{El reglamento de quejas universitarias}

Creamos un documento que cubre:

\begin{itemize}
  \item Cap\'itulo 1: Disposiciones generales (5 art\'iculos)
  \item Cap\'itulo 2: Proceso de atenci\'on (4 art\'iculos)
  \item Cap\'itulo 3: Derechos del reportante (3 art\'iculos)
  \item Cap\'itulo 4: Sanciones (3 art\'iculos)
  \item Cap\'itulo 5: Procedimiento de escalamiento (3 art\'iculos)
  \item FAQ de acoso
  \item FAQ de infraestructura
  \item Contacto
\end{itemize}

\begin{lstlisting}[language=bash,caption={knowledge/reglamento.txt}]
REGLAMENTO DE QUEJAS UNIVERSITARIAS - SENTINEL ACADEMIA
======================================================

Capitulo 1: Disposiciones generales

Articulo 1. Toda queja presentada por un estudiante o docente
de la universidad sera registrada en el sistema Sentinel AcademIA
para su seguimiento y resolucion.

Articulo 2. Las quejas pueden ser anonimas o con identificacion
del reportante. Las quejas anonimas reciben el mismo tratamiento
que las identificadas.

Articulo 3. Toda queja debe incluir una descripcion detallada
del problema con un minimo de 20 caracteres. Se recomienda
adjuntar evidencia fotografica o documental cuando sea posible.

Articulo 4. Las categorias de quejas son:
- ACADEMICA: problemas con profesores, cursos, calificaciones
- INFRAESTRUCTURA: problemas con aulas, edificios, equipos
- ACOSO: situaciones de acoso, hostigamiento, violencia
- ADMINISTRATIVA: problemas con matriculas, tramites
- SALUD: temas de salud mental, fisica
- OTRA: cualquier otra situacion
# ... 4KB total
\end{lstlisting}

\subsection{Subir el documento a OCI}

\begin{lstlisting}[language=Python,caption={Upload to OCI}]
object_storage.put_object(
    namespace,
    bucket,
    "reglamento.txt",
    data,  # bytes del archivo
    content_type="text/plain",
    content_encoding="utf-8",
)
\end{lstlisting}

% ============================================================
\section{Paso 3: knowledge\_service.py (retrieval)}
% =========================================================%

\subsection{El patr\'on: keyword-based search}

Para evitar el package \texttt{oci} (229MB), hacemos retrieval en Python puro.

\begin{lstlisting}[language=Python,caption={src/services/knowledge\_service.py (core)}]
import re

KNOWLEDGE_BASE: list[dict[str, str]] = [
    {
        "id": "reglamento-cap1",
        "title": "Capitulo 1: Disposiciones generales",
        "content": (
            "Articulo 1. Toda queja sera registrada en el sistema "
            "Sentinel AcademIA. Articulo 4. Las categorias son: "
            "ACADEMICA (profesores, cursos), INFRAESTRUCTURA (aulas), "
            "ACOSO (hostigamiento), ADMINISTRATIVA (matriculas), "
            "SALUD (salud mental), OTRA."
        ),
    },
    # ... 7 documents mas
]


def _tokenize(text: str) -> list[str]:
    """Tokenize text into lowercase words (no punctuation)."""
    return re.findall(r"\b\w+\b", text.lower())


def _score_doc(query_tokens: list[str], doc: dict[str, str]) -> int:
    """Score a document by how many query tokens it contains."""
    doc_tokens = set(_tokenize(doc["title"] + " " + doc["content"]))
    return sum(1 for tok in query_tokens if tok in doc_tokens)


def search(query: str, top_k: int = 3) -> list[dict]:
    """Search the knowledge base for the most relevant documents."""
    if not query or not query.strip():
        return []

    query_tokens = _tokenize(query)
    if not query_tokens:
        return []

    scored = [(doc, _score_doc(query_tokens, doc)) for doc in KNOWLEDGE_BASE]
    scored = [(doc, score) for doc, score in scored if score > 0]
    scored.sort(key=lambda x: x[1], reverse=True)

    return [
        {
            "id": doc["id"],
            "title": doc["title"],
            "content": doc["content"],
            "score": score,
        }
        for doc, score in scored[:top_k]
    ]


def build_context(query: str, top_k: int = 3) -> str:
    """Build a context string for the LLM from the search results."""
    results = search(query, top_k=top_k)
    if not results:
        return ""

    lines = ["[Contexto relevante del reglamento:]"]
    for r in results:
        lines.append(f"\n--- {r['title']} ---")
        lines.append(r["content"])
    return "\n".join(lines)
\end{lstlisting}

\textbf{Decisiones}:
\begin{itemize}
  \item \ok \textbf{8 documentos hardcoded} (en lugar de cargar de OCI en runtime)
  \item \ok \textbf{Keyword-based} (en lugar de embeddings)
  \item \ok \textbf{Top-K=3}: 3 documentos m\'as relevantes
  \item \ok \textbf{Para producci\'on}: reemplazar con embeddings + vector store
\end{itemize}

% ============================================================
\section{Paso 4: Schemas (ChatRequest / ChatResponse)}
% ====================================

\subsection{Los schemas alineados con OpenAPI}

\begin{lstlisting}[language=Python,caption={src/handlers/\_shared/schemas.py}]
class ChatRequest(BaseModel):
    """POST /api/chat - User question for the RAG-powered chatbot."""
    question: str = Field(min_length=5, max_length=1000)
    context: str | None = Field(default="TODOS")
    conversationId: str | None = None


class ChatSource(BaseModel):
    """Source document referenced in the answer."""
    id: str
    title: str
    content: str
    score: float = Field(ge=0, le=1)


class ChatResponse(BaseModel):
    """Response from /api/chat."""
    answer: str
    sources: list[ChatSource] = Field(default_factory=list)
    conversationId: str | None = None
    modeloUsado: str
    tokensUsados: int = 0
    latenciaMs: int = 0
\end{lstlisting}

% ============================================================
\section{Paso 5: Lambda chat.py}
% ====================================

\subsection{El handler completo}

\begin{lstlisting}[language=Python,caption={src/handlers/chat.py (handler)}]
SYSTEM_PROMPT = """Eres un asistente virtual de la universidad Sentinel AcademIA.
Tu trabajo es responder preguntas de estudiantes y personal sobre el reglamento
de quejas universitarias, basandote SIEMPRE en el contexto proporcionado.

Reglas:
1. Responde en espanol, de forma clara y concisa.
2. Basa tus respuestas EXCLUSIVAMENTE en el contexto del reglamento.
3. Si el contexto no contiene la respuesta, dilo claramente.
4. Cita el articulo o capitulo relevante cuando sea posible.
5. Si la pregunta es sobre una situacion personal, sugiere escalar la queja.
6. NO inventes informacion que no este en el contexto."""


def _generate_answer_mock(question: str, context: str) -> str:
    """Mock answer when LLM is not available."""
    if not context:
        return (
            "No encontre informacion relevante en el reglamento para tu pregunta. "
            "Puedes consultar el capitulo 1 o contactar a quejas@universidad.edu."
        )
    return (
        f"Basado en el reglamento, los articulos mas relevantes para tu "
        f"pregunta son:\n\n{context[:600]}...\n\n"
        f"Para mas detalles, te recomiendo consultar el documento completo."
    )


def _generate_answer_with_llm(question: str, context: str):
    """Generate an answer using the LLM (with context)."""
    import os
    use_mock = os.environ.get("MOCK_LLM", "true").lower() != "false"
    start = time.time()
    if use_mock:
        answer = _generate_answer_mock(question, context)
        tokens = len(question.split()) + len(context.split())
    else:
        # Real Cohere call would go here
        ...
    latency_ms = int((time.time() - start) * 1000)
    return answer, tokens, latency_ms


@error_handler
def lambda_handler(event, context):
    with_correlation_id(event)
    try:
        tenant_id = extract_tenant_id(event)
        log.append_keys(tenant_id=tenant_id)

        body = _parse_body(event)
        req = ChatRequest.model_validate(body)

        # 1. Retrieve relevant context (RAG step 1)
        context = build_context(req.question, top_k=3)
        sources_raw = search(req.question, top_k=3)

        # 2. Generate answer with context (RAG step 2)
        answer, tokens, latency = _generate_answer_with_llm(req.question, context)

        # 3. Build response
        sources = [
            ChatSource(
                id=s["id"],
                title=s["title"],
                content=s["content"],
                score=min(1.0, s["score"] / 10.0),
            )
            for s in sources_raw
        ]

        return HttpResponse.ok(ChatResponse(
            answer=answer,
            sources=sources,
            conversationId=req.conversationId,
            modeloUsado="cohere.command-latest",
            tokensUsados=tokens,
            latenciaMs=latency,
        ).model_dump(mode="json"))
    finally:
        clear_correlation_id()
\end{lstlisting}

% ============================================================
\section{Paso 6: SAM template (la nueva Lambda)}
% ====================================

\subsection{ChatFunction}

\begin{lstlisting}[language=yaml,caption={infra/template.yaml — ChatFunction}]
ChatFunction:
  Type: AWS::Serverless::Function
  Properties:
    Role: !Ref LambdaRoleArn
    FunctionName: !Sub "sentinel-chat-${Stage}"
    CodeUri: ../src
    Handler: handlers.chat.lambda_handler
    MemorySize: 512
    Timeout: 30
    Environment:
      Variables:
        MOCK_LLM: "true"  # <-- usa MockLLMClient (no oci package)
    Events:
      ChatApi:
        Type: Api
        Properties:
          Path: /api/chat
          Method: POST
          RestApiId: !Ref SentinelApi
\end{lstlisting}

\textbf{Decisiones}:
\begin{itemize}
  \item \ok \textbf{No permisos S3 ni DynamoDB}: el chat no toca storage
  \item \ok \textbf{MOCK\_LLM=true}: el Lambda no tiene \texttt{oci} package
  \item \ok \textbf{Timeout 30s}: el LLM real puede tardar (Cohere ~2-5s)
\end{itemize}

% ============================================================
\section{El bug que arreglamos}
% ====================================

\begin{tcolorbox}[colback=errred!5!white,colframe=errred,breakable,title={Bug: KeyError 'content'}]
\textbf{Cuando}: el primer test del chat retornaba 500. \\
\textbf{Causa}: usaba \texttt{get\_top\_sources} que NO inclu\'ia content, pero en el loop acced\'ia \texttt{s["content"]}.

\begin{lstlisting}
# MAL
sources_raw = get_top_sources(query, top_k=3)
# ... despues ...
content=s["content"]  # <-- KeyError!
\end{lstlisting}

\textbf{Soluci\'on}: usar \texttt{search} directamente (que s\'i tiene content).

\begin{lstlisting}
# BIEN
sources_raw = search(query, top_k=3)
# ahora sources_raw tiene 'content' en cada item
\end{lstlisting}
\end{tcolorbox}

% ============================================================
\section{Los 2 tests E2E que validan el RAG}
% ====================================

\subsection{Test 1: plazos de queja cr\'itica}

\begin{lstlisting}[language=bash,caption={Pregunta sobre plazos}]
curl -s -X POST "$API/api/chat" \
  -H "Content-Type: application/json" \
  -H "X-Tenant-ID: demo-utpl" \
  -d '{"question":"Cuanto tiempo tengo para que se resuelva una queja critica?"}'
\end{lstlisting}

\begin{lstlisting}[language=json,caption={Output esperado}]
{
  "answer": "Basado en el reglamento, los articulos mas relevantes para tu pregunta son: ...",
  "sources": [
    {
      "id": "reglamento-faq-acoso",
      "title": "FAQ: Que hacer en caso de acoso",
      "content": "Si sufres acoso: 1) Reporta inmediatamente...",
      "score": 0.5
    },
    {
      "id": "reglamento-faq-infraestructura",
      "title": "FAQ: Problemas de infraestructura",
      "content": "Para aulas danadas...",
      "score": 0.2
    },
    {
      "id": "reglamento-cap1",
      "title": "Capitulo 1: Disposiciones generales",
      "content": "Articulo 1. Toda queja sera registrada...",
      "score": 0.1
    }
  ],
  "tokensUsados": 146,
  "latenciaMs": 0
}
\end{lstlisting}

\subsection{Test 2: situaci\'on de acoso}

\begin{lstlisting}[language=bash,caption={Pregunta sobre acoso}]
curl -s -X POST "$API/api/chat" \
  -H "Content-Type: application/json" \
  -H "X-Tenant-ID: demo-utpl" \
  -d '{"question":"Que hago si sufro acoso en la universidad?"}'
\end{lstlisting}

\begin{lstlisting}[language=json,caption={Output esperado}]
{
  "answer": "...",
  "sources": [
    { "id": "reglamento-faq-acoso", "score": 0.4 },
    { "id": "reglamento-cap4", "score": 0.3 },
    { "id": "reglamento-faq-infraestructura", "score": 0.3 }
  ]
}
\end{lstlisting}

% ============================================================
\section{C\'omo escalar a OCI Generative AI Agents (producci\'on)}
% ====================================

\subsection{¿Qu\'e es OCI Generative AI Agents?}

Es el servicio managed de Oracle para RAG. Oracle hace TODO el trabajo pesado.

\subsection{Setup paso a paso}

\subsubsection{Paso 1: Crear Object Storage bucket}

\begin{lstlisting}[language=bash,caption={Ya lo hicimos}]
oci os bucket create \
  --compartment-id $COMPARTMENT \
  --name sentinel-knowledge
\end{lstlisting}

\subsubsection{Paso 2: Subir los documentos}

\begin{lstlisting}[language=bash,caption={Ya lo hicimos}]
oci os object put \
  --bucket-name sentinel-knowledge \
  --name reglamento.txt \
  --file knowledge/reglamento.txt
\end{lstlisting}

\subsubsection{Paso 3: Crear Data Source (en Consola OCI)}

\begin{enumerate}
  \item Ir a la Consola de OCI
  \item Men\'u $\rightarrow$ \textbf{Analytics \& AI} $\rightarrow$ \textbf{Generative AI Agents}
  \item Click \textbf{Create Agent}
  \item Seleccionar el compartment \texttt{SentinelAcademia}
  \item Step 1: \textbf{Data source}: Object Storage, bucket \texttt{sentinel-knowledge}
  \item Step 2: Configurar chunking (default: 500 tokens, 50 overlap)
  \item Step 3: Crear index (vector embeddings con Cohere embed-multilingual-v3.0)
  \item Esperar 5-10 minutos a que indexe
\end{enumerate}

\subsubsection{Paso 4: Crear el Agent}

\begin{enumerate}
  \item En el wizard: \textbf{Inference model}: Cohere Command R
  \item \textbf{Instructions}: "Responde bas\'andote en el contexto. Cita el art\'iculo."
  \item \textbf{Enable citations}: true
  \item Click \textbf{Create}
\end{enumerate}

\subsubsection{Paso 5: Crear Endpoint}

\begin{enumerate}
  \item Click en el Agent
  \item Tab \textbf{Endpoints} $\rightarrow$ \textbf{Create endpoint}
  \item Obtener el \textbf{OCID} del endpoint
  \item URL: \texttt{https://agent-runtime.generativeai.us-chicago-1.oci.oraclecloud.com}
\end{enumerate}

\subsubsection{Paso 6: Llamar desde Python (en Lambda)}

\begin{lstlisting}[language=Python,caption={oci\_agent.py — Cliente de OCI Agent}]
import oci
from oci.config import from_file

config = from_file("/home/lFunknown/.oci/private_key.pem", "DEFAULT")
# Nota: 229MB, no cabe en Lambda
# Pero para DESKTOP/LOCAL es OK

agent_runtime_client = oci.generative_ai_agent_runtime.GenerativeAiAgentRuntimeClient(
    config,
    service_endpoint="https://agent-runtime.generativeai.us-chicago-1.oci.oraclecloud.com"
)

response = agent_runtime_client.chat(
    agent_endpoint_id="ocid1.genaiagentendpoint.oc1.us-chicago-1.aaa...",
    chat_details=oci.generative_ai_agent_runtime.models.ChatDetails(
        user_message="Que hago si sufro acoso?",
        session_id="session-123",
    )
)

print(response.data.message)
print("Citations:", response.data.citations)
\end{lstlisting}

\subsection{Para usar OCI Agent en Lambda}

\textbf{Workarounds}:
\begin{enumerate}
  \item \textbf{Lambda Container Image}: empaqueta tu propia imagen Docker con el SDK
  \item \textbf{Lambda Layer}: sube el SDK \texttt{oci} como un layer (\textasciitilde 50MB slim)
  \item \textbf{Llamar via HTTPS}: usa \texttt{requests} con signing manual de OCI
  \item \textbf{API Gateway proxy}: tu Lambda hace fetch al endpoint del Agent
\end{enumerate}

% ============================================================
\section{Cat\'alogo de errores y soluciones}
% ====================================

\begin{tcolorbox}[colback=errred!5!white,colframe=errred,breakable,title={Error 1: Namespace mismatch}]
\textbf{Cuando}: \texttt{RelatedResourceNotAuthorizedOrNotFound}. \\
\textbf{Causa}: pasamos el OCID del tenancy como namespace. \\
\textbf{Soluci\'on}: usar el namespace real (\texttt{ax59x2so1pxw}), NO el OCID.
\end{tcolorbox}

\begin{tcolorbox}[colback=errred!5!white,colframe=errred,breakable,title={Error 2: KeyError 'content' en chat}]
\textbf{Cuando}: 500 al primer test del chat. \\
\textbf{Causa}: \texttt{get\_top\_sources} no inclu\'ia content. \\
\textbf{Soluci\'on}: usar \texttt{search} directamente.
\end{tcolorbox}

\begin{tcolorbox}[colback=errred!5!white,colframe=errred,breakable,title={Error 3: 502 al deployar con oci package}]
\textbf{Cuando}: \texttt{sam deploy} con \texttt{oci} en requirements. \\
\textbf{Causa}: package de 229MB excede 250MB de Lambda. \\
\textbf{Soluci\'on}: MockLLMClient en Lambda, RealLLMClient solo local.
\end{tcolorbox}

\begin{tcolorbox}[colback=errred!5!white,colframe=errred,breakable,title={Error 4: chat endpoint 404}]
\textbf{Cuando}: "Missing Authentication Token" en POST /api/chat. \\
\textbf{Causa}: el API Gateway no tiene la ruta (no se re-despleg\'o). \\
\textbf{Soluci\'on}: \texttt{sam deploy} para actualizar el API.
\end{tcolorbox}

\begin{tcolorbox}[colback=errred!5!white,colframe=errred,breakable,title={Error 5: No encuentra fuentes relevantes}]
\textbf{Cuando}: \texttt{sources: []} y answer gen\'erico. \\
\textbf{Causa}: pregunta sin overlap con knowledge base. \\
\textbf{Soluci\'on}: expandir knowledge base o usar embeddings sem\'anticos.
\end{tcolorbox}

% ============================================================
\section{Lo que aprendimos (resumen ejecutivo)}
% ====================================

\begin{tcolorbox}[colback=okgreen!5!white,colframe=okgreen,title={\textbf{Checklist Fase 5}}]
\begin{enumerate}
  \item \textbf{RAG = retrieval + generation}: retrieve relevant docs, then LLM
  \item \textbf{OCI Object Storage} = S3 de Oracle (compatible con boto3)
  \item \textbf{Namespace} vs OCID: el namespace es string corto, NO el OCID
  \item \textbf{oci package} = 229MB, no cabe en Lambda $\rightarrow$ MockLLM
  \item \textbf{Keyword search} en Lambda (sin deps externas) para el demo
  \item \textbf{Top-K=3} docs m\'as relevantes para el LLM
  \item \textbf{Citations} en la respuesta (id, title, content, score)
  \item \textbf{System prompt} estricto: "solo responde con el contexto"
  \item \textbf{OCI Agent} (producci\'on) maneja TODO el pipeline de RAG
  \item \textbf{Lambda Layer} o \textbf{Container Image} para usar oci en Lambda
\end{enumerate}
\end{tcolorbox}

% ============================================================
\section{Ejercicios para practicar}
% ====================================

\subsection{Ejercicio 1: Embeddings sem\'anticos}

Reemplazar keyword search con embeddings usando \texttt{cohere.embed-multilingual-v3.0}. Calcular cosine similarity. Subir a FAISS o usar OCI Search.

\subsection{Ejercicio 2: Multi-turn conversations}

A\~nadir \texttt{conversationId} y guardar el historial en DynamoDB. El LLM recibe historial + pregunta.

\subsection{Ejercicio 3: Streaming de respuestas}

Usar Server-Sent Events (SSE) para que el LLM responda token por token.

\subsection{Ejercicio 4: OCI Agent real}

Desplegar OCI Generative AI Agent en la Consola y conectarlo desde una Lambda con Layer.

\subsection{Ejercicio 5: Hybrid search}

Combinar keyword search (r\'apido) con embeddings (preciso). Re-rank con Cohere.

% ============================================================
\section{Siguiente paso: Fase 6 (Testing)}
% ====================================

En la \textbf{Fase 6} a\~nadimos tests:

\begin{itemize}
  \item \texttt{pytest} con \texttt{moto} (mock de AWS services)
  \item Tests unitarios de cada Lambda
  \item Tests de integraci\'on E2E
  \item Coverage $>$ 70\%
\end{itemize}

\textbf{Tiempo estimado}: 1-1.5 horas.

\begin{center}
\begin{tabular}{ll}
\toprule
Fase & Estado \\
\midrule
1. Backend Core & [OK] \\
2. LLM integration & [OK] \\
3. Frontend deploy & [OK] \\
4. Files & [OK] \\
5. RAG (este tutorial) & [OK] \\
6. Testing & pr\'oximo \\
7. Deploy final + demo & pendiente \\
\bottomrule
\end{tabular}
\end{center}

\vspace{2em}
\begin{center}
\textit{Fin del tutorial. \`A por la Fase 6!}
\end{center}

\end{document}
